\section{Methodology}
\label{sec:method}

We study the dimensionality of belief space through three complementary PCA analyses, each capturing a different aspect of belief structure. \figref{fig:comparison} provides an overview of how the three methods compare.

\subsection{Belief Generation}
\label{sec:belief_gen}

We constructed a corpus of \nbeliefs unique belief statements from two sources. First, we drew 1,226 existing belief questions from \opinionqa~\cite{santurkar2023whose} and curated datasets, yielding 2,452 perspective statements (each question produces both an affirmative and a negating stance). Second, we generated 7,541 new beliefs using \gptmini across 50 topic categories (politics, economics, religion, science, education, healthcare, environment, social justice, family, gender, immigration, criminal justice, and 38 others) combined with 20 angle modifiers to ensure diversity within each topic. After deduplication, the final corpus contains \nbeliefs unique belief statements.

\subsection{Experiment 1: Semantic Embedding PCA}
\label{sec:embed_pca}

We embedded all \nbeliefs beliefs using OpenAI's \texttt{text-embedding-3-small} model (1,536 dimensions). After standardizing the embedding matrix, we applied PCA to extract the principal components of semantic belief space. This analysis captures {\em topical} structure---which beliefs discuss similar subjects---independent of whether people who hold one belief tend to hold another.

\subsection{Experiment 2: Persona Response PCA}
\label{sec:response_pca}

\para{Belief selection.}
To make the response experiment tractable, we selected a diverse subset of \nbeliefsub beliefs from the full corpus. We applied $k$-means clustering ($k=20$) in embedding space and selected 10 beliefs closest to each cluster center plus 5 random beliefs per cluster, ensuring maximal topical coverage.

\para{Persona generation.}
We generated \npersonas synthetic personas by sampling from a combinatorial demographic space: 7 age groups $\times$ 3 genders $\times$ 5 education levels $\times$ 10 political orientations $\times$ 10 religions $\times$ 15 countries $\times$ 8 personality types. Each persona is defined by a unique combination of these attributes.

\para{Response collection.}
Each persona rated all \nbeliefsub beliefs on a 5-point Likert scale (1 = Strongly Disagree to 5 = Strongly Agree) using \gptmini with temperature 0.7. This produced a $300 \times 300$ response matrix $\mR$ with no missing values.

\para{PCA procedure.}
We standardized columns of $\mR$ (zero mean, unit variance per belief) and applied PCA. We compared the resulting eigenvalue spectrum against a random baseline generated by shuffling each column of $\mR$ independently, destroying covariance structure while preserving marginal distributions.

\subsection{Experiment 3: Conditional Belief Influence PCA}
\label{sec:conditional_pca}

To measure the {\em causal} structure of beliefs---how knowing someone's stance on one belief predicts their stances on others---we designed a conditional influence experiment. We selected \nanchors anchor beliefs and \ntestbeliefs test beliefs. For each anchor belief, we prompted \gptmini (temperature 0.3) with: ``Someone who strongly agrees with [anchor]---how would they rate [test beliefs]?'' and the symmetric query for ``strongly disagrees.'' We computed the difference matrix $\mD \in \mathbb{R}^{50 \times 100}$, where $D_{ij}$ is the change in predicted rating of test belief $j$ when the anchor belief $i$ is agreed vs.\ disagreed with. PCA on $\mD$ reveals the latent dimensions of belief influence.

\subsection{Evaluation Metrics}

We report the following metrics for each PCA analysis:
\begin{itemize}[leftmargin=*,itemsep=0pt,topsep=0pt]
    \item {\bf Components for threshold:} number of principal components needed to explain 50\%, 80\%, 90\%, and 95\% of total variance.
    \item {\bf Kaiser criterion:} number of components with eigenvalue $> 1$ (after standardization).
    \item {\bf Scree analysis:} visual inspection of eigenvalue decay.
    \item {\bf Component interpretability:} top-loading beliefs on each principal component, assessed qualitatively.
\end{itemize}

\para{Baselines.}
We compare against two baselines: (1)~a {\em random baseline} constructed by independently shuffling each column of $\mR$, which preserves marginal distributions but destroys covariance; and (2)~the {\em theoretical uniform} case where $n$ independent dimensions produce eigenvalues $\lambda_i = 1$ for all $i$.

\subsection{Reproducibility}

All experiments use random seed 42. Response collection used \gptmini with temperature 0.7 (persona responses) or 0.3 (conditional queries). The response matrix, embeddings, and all intermediate results are saved as NumPy arrays. The full pipeline requires approximately 2,700 API calls.
